% ===== PRÉAMBULE CANONIQUE — à coller VERBATIM en tête de CHAQUE .tex =====
\documentclass[11pt,a4paper]{article}
\usepackage[utf8]{inputenc}\usepackage[T1]{fontenc}\usepackage{lmodern}
\usepackage[french]{babel}
\usepackage{amsmath,amssymb,mathtools}\usepackage{icomma}
\usepackage[version=4]{mhchem}          % \ce{...} : équations & formules
\usepackage{siunitx}\sisetup{locale=FR} % \qty{}{}, \si{}, \ang{}
\usepackage{chemfig}                    % \chemfig{...} : structures développées
\usepackage{xcolor}\usepackage{colortbl}
\usepackage{tikz}\usepackage{pgfplots}\pgfplotsset{compat=1.18}
\usepackage[most]{tcolorbox}\usepackage{enumitem}\usepackage{array,booktabs}
\usepackage[a4paper,margin=2.2cm]{geometry}
\usetikzlibrary{arrows.meta,calc,angles,quotes,positioning,patterns,backgrounds,shapes.geometric}
\definecolor{primary}{RGB}{222,49,99}\definecolor{primarydark}{RGB}{150,22,55}
\definecolor{defcol}{RGB}{0,114,178}\definecolor{methcol}{RGB}{52,92,148}
\definecolor{excol}{RGB}{0,135,90}\definecolor{attncol}{RGB}{200,40,55}
\tcbset{boxbase/.style={enhanced,breakable,boxrule=0.9pt,arc=2.5pt,
  left=3mm,right=3mm,top=2.3mm,bottom=2.3mm,fonttitle=\bfseries,coltitle=white}}
% --- boîtes TUTORIEL (noms EXACTS, ne pas renommer) ---
\newtcolorbox{cle}[1][]{boxbase,colback=defcol!6,colframe=defcol,
  title={\textsc{À retenir}\ifx\relax#1\relax\relax\else\textmd{ : \itshape #1}\fi}}
\newtcolorbox{methode}[1][]{boxbase,colback=methcol!7,colframe=methcol,sharp corners,
  title={\textsc{Méthode}\ifx\relax#1\relax\relax\else\textmd{ --- \itshape #1}\fi}}
\newtcolorbox{exemplebox}[1][]{boxbase,colback=excol!5,colframe=excol,
  title={\textsc{Exemple}\ifx\relax#1\relax\relax\else\textmd{ : \itshape #1}\fi}}
\newtcolorbox{piege}[1][]{boxbase,colback=attncol!6,colframe=attncol,
  title={\textsc{Piège}\ifx\relax#1\relax\relax\else\textmd{ : \itshape #1}\fi}}
\newtcolorbox{remarque}[1][]{boxbase,colback=black!4,colframe=black!45,
  title={\textsc{Remarque}\ifx\relax#1\relax\relax\else\textmd{ : \itshape #1}\fi}}
% --- boîtes EXERCICE / CORRIGÉ (noms EXACTS, auto-comptées) ---
\newtcolorbox[auto counter]{exo}[1][]{boxbase,colback=primary!4,colframe=primary,
  title={\textsc{Exercice}~\thetcbcounter\ifx\relax#1\relax\relax\else\textmd{ --- \itshape #1}\fi}}
\newtcolorbox[auto counter]{corrige}[1][]{boxbase,colback=excol!4,colframe=excol,
  title={\textsc{Corrigé}~\thetcbcounter\ifx\relax#1\relax\relax\else\textmd{ --- \itshape #1}\fi}}
\newcommand{\R}{\mathbb{R}}\newcommand{\N}{\mathbb{N}}\newcommand{\Z}{\mathbb{Z}}\newcommand{\ds}{\displaystyle}
% (un \usepackage{fancyhdr}+en-tête est possible ; il est ignoré côté web)
% =========================================================================
\begin{document}
\sisetup{per-mode=symbol}

\begin{exo}[de la concentration au pH]
Calcule le pH de chaque solution ($\mathrm{pH} = -\log[\ce{H3O+}]$).
\begin{enumerate}[label=\alph*)]
  \item $[\ce{H3O+}] = \qty{e-3}{mol\per L}$
  \item $[\ce{H3O+}] = \qty{e-11}{mol\per L}$
  \item $[\ce{H3O+}] = \qty{1,0e-5}{mol\per L}$
  \item $[\ce{H3O+}] = \qty{0,010}{mol\per L}$
\end{enumerate}
\end{exo}

\begin{exo}[du pH à la concentration]
Calcule $[\ce{H3O+}]$ de chaque solution ($[\ce{H3O+}] = 10^{-\mathrm{pH}}$).
\begin{enumerate}[label=\alph*)]
  \item $\mathrm{pH} = 2$
  \item $\mathrm{pH} = 9$
  \item $\mathrm{pH} = 7$
  \item $\mathrm{pH} = 1$
\end{enumerate}
\end{exo}

\begin{exo}[le produit ionique de l'eau]
À \qty{25}{\celsius}, $K_e = [\ce{H3O+}][\ce{OH-}] = 10^{-14}$. Complète.
\begin{enumerate}[label=\alph*)]
  \item $[\ce{H3O+}] = \qty{e-4}{mol\per L}$ : calcule $[\ce{OH-}]$.
  \item $[\ce{OH-}] = \qty{e-2}{mol\per L}$ : calcule $[\ce{H3O+}]$.
  \item $\mathrm{pH} = 4$ : calcule le pOH.
\end{enumerate}
\end{exo}

\begin{exo}[acide, neutre ou basique ? (à \qty{25}{\celsius})]
Indique le caractère de chaque solution.
\begin{enumerate}[label=\alph*)]
  \item $\mathrm{pH} = 3$
  \item $\mathrm{pH} = 7$
  \item $\mathrm{pH} = 11$
  \item $[\ce{H3O+}] = \qty{e-9}{mol\per L}$
\end{enumerate}
\end{exo}

\begin{exo}[de $K_a$ au $\mathrm{p}K_a$, et classement des forces]
\begin{enumerate}[label=\alph*)]
  \item Calcule le $\mathrm{p}K_a$ d'un acide de constante $K_a = \qty{e-5}{}$.
  \item Voici quatre acides : \ce{HF} ($\mathrm{p}K_a = 3{,}2$), \ce{CH3COOH} ($4{,}8$), \ce{HClO} ($7{,}5$),
        \ce{NH4+} ($9{,}2$). Range-les du \textbf{plus fort} au \textbf{plus faible}.
\end{enumerate}
\end{exo}

\end{document}
